\documentclass[10pt,a4paper]{amsart}
\usepackage{amsmath,amssymb,amsthm,graphicx}
\usepackage{verbatim}
\newenvironment{reference}{\comment}{\endcomment}
%\newenvironment{reference}{}{}
\newenvironment{slogan}{\comment}{\endcomment}
\newenvironment{history}{\comment}{\endcomment}

% For commutative diagrams we use Xy-pic
\usepackage[all]{xy}

% We use 2cell for 2-commutative diagrams.
\xyoption{2cell}
% Theorem environments.
%
\theoremstyle{plain}
\newtheorem{theorem}[subsection]{Theorem}
\newtheorem{proposition}[subsection]{Proposition}
\newtheorem{lemma}[subsection]{Lemma}

\theoremstyle{definition}
\newtheorem{definition}[subsection]{Definition}
\newtheorem{example}[subsection]{Example}
\newtheorem{exercise}[subsection]{Exercise}
\newtheorem{situation}[subsection]{Situation}

\theoremstyle{remark}
\newtheorem{remark}[subsection]{Remark}
\newtheorem{remarks}[subsection]{Remarks}

\numberwithin{equation}{subsection}

% Macros
%
\def\lim{\mathop{\mathrm{lim}}\nolimits}
\def\colim{\mathop{\mathrm{colim}}\nolimits}
\def\Spec{\mathop{\mathrm{Spec}}}
\def\Hom{\mathop{\mathrm{Hom}}\nolimits}
\def\Ext{\mathop{\mathrm{Ext}}\nolimits}
\def\SheafHom{\mathop{\mathcal{H}\!\mathit{om}}\nolimits}
\def\SheafExt{\mathop{\mathcal{E}\!\mathit{xt}}\nolimits}
\def\Sch{\mathit{Sch}}
\def\Mor{\mathop{\mathrm{Mor}}\nolimits}
\def\Ob{\mathop{\mathrm{Ob}}\nolimits}
\def\Sh{\mathop{\mathit{Sh}}\nolimits}
\def\NL{\mathop{N\!L}\nolimits}
\def\CH{\mathop{\mathrm{CH}}\nolimits}
\def\proetale{{pro\text{-}\acute{e}tale}}
\def\etale{{\acute{e}tale}}
\def\QCoh{\mathit{QCoh}}
\def\Ker{\mathop{\mathrm{Ker}}}
\def\Im{\mathop{\mathrm{Im}}}
\def\Coker{\mathop{\mathrm{Coker}}}
\def\Coim{\mathop{\mathrm{Coim}}}

% Boxtimes
%
\DeclareMathSymbol{\boxtimes}{\mathbin}{AMSa}{"02}

%
% Macros for moduli stacks/spaces
%
\def\QCohstack{\mathcal{QC}\!\mathit{oh}}
\def\Cohstack{\mathcal{C}\!\mathit{oh}}
\def\Spacesstack{\mathcal{S}\!\mathit{paces}}
\def\Quotfunctor{\mathrm{Quot}}
\def\Hilbfunctor{\mathrm{Hilb}}
\def\Curvesstack{\mathcal{C}\!\mathit{urves}}
\def\Polarizedstack{\mathcal{P}\!\mathit{olarized}}
\def\Complexesstack{\mathcal{C}\!\mathit{omplexes}}
% \Pic is the operator that assigns to X its picard group, usage \Pic(X)
% \Picardstack_{X/B} denotes the Picard stack of X over B
% \Picardfunctor_{X/B} denotes the Picard functor of X over B
\def\Pic{\mathop{\mathrm{Pic}}\nolimits}
\def\Picardstack{\mathcal{P}\!\mathit{ic}}
\def\Picardfunctor{\mathrm{Pic}}
\def\Deformationcategory{\mathcal{D}\!\mathit{ef}}
\usepackage{cleveref}
\crefname{originaldefinition}{original definition}{original definitions}
\Crefname{originaldefinition}{Original definition}{Original definitions}
\makeatletter
\newlabel{formal-spaces-definition-adic-homomorphism}{{0GBR}{source}{}{sourcepointer-1}{}}
\newlabel{formal-spaces-definition-adic-homomorphism@cref}{{[originaldefinition][0][]0GBR}{[1][0][]source}{}{}{}}
\makeatother
\crefname{originalremark}{original remark}{original remarks}
\Crefname{originalremark}{Original remark}{Original remarks}
\makeatletter
\newlabel{formal-spaces-remark-base-change-variant-variant-Noetherian}{{0GBG}{source}{}{sourcepointer-2}{}}
\newlabel{formal-spaces-remark-base-change-variant-variant-Noetherian@cref}{{[originalremark][0][]0GBG}{[1][0][]source}{}{}{}}
\makeatother
\crefname{originalreference}{original reference}{original references}
\Crefname{originalreference}{Original reference}{Original references}
\makeatletter
\newlabel{lemma-base-change-rig-smooth-homomorphism}{{0GAM}{source}{}{sourcepointer-3}{}}
\newlabel{lemma-base-change-rig-smooth-homomorphism@cref}{{[originalreference][0][]0GAM}{[1][0][]source}{}{}{}}
\makeatother
\crefname{originalreference}{original reference}{original references}
\Crefname{originalreference}{Original reference}{Original references}
\makeatletter
\newlabel{lemma-rig-smooth-continuous}{{0GCI}{source}{}{sourcepointer-4}{}}
\newlabel{lemma-rig-smooth-continuous@cref}{{[originalreference][0][]0GCI}{[1][0][]source}{}{}{}}
\makeatother
\title{Rig-smoothness is stable under adic base change in WAdm\textasciicircum{}Noeth}
\author{The Stacks Project Authors}
\date{stacks-project:restricted.tex:local-snapshot-20261012; GNU FDL 1.2 or later}
\begin{document}
\maketitle
\noindent\textit{Editorial note.} All delivered mathematical text is the author's own, copied from the pinned Stacks Project chapter file restricted.tex; no statement or proof step was rewritten, reordered or removed.\par\medskip
\noindent\textit{Editorial note.} The counted proof is the author's complete argument: the proof chooses a common ideal of definition and delegates the rig-smoothness of the completed tensor product to the homomorphism-level base-change lemma\par\medskip
\noindent\textit{Editorial note.} The label lemma-base-change-rig-smooth-continuous names the counted statement, and the proof environment immediately following it in the source is delivered verbatim as its proof.\par\medskip
\noindent\textit{Editorial note.} The counted statement is carried by the author's own lemma environment in the chapter source.\par\medskip
\noindent\textit{Allowed standard prerequisites (editorial).}\begin{itemize}
\item rig-smooth continuous families and adic morphisms in the category WAdm\textasciicircum{}Noeth (lemma-rig-smooth-continuous and Formal Spaces definition of adic homomorphism)
\item base change of rig-smooth homomorphisms of adic rings (lemma-base-change-rig-smooth-homomorphism of the same chapter)
\end{itemize}\medskip
\ifdefined\hypertarget\hypertarget{sourcepointer-1}{}\fi
\noindent\textit{Source reference.} Original definition 0GBR: sources/stacks-chow/tags.txt, source lines 20381--20381. Prerequisite pointer invoked at proof line(s) 4385; it names a standard fact of a sibling chapter that the counted argument uses, and every step of the counted proof built on it is retained verbatim in the delivered body.\par\medskip
\ifdefined\hypertarget\hypertarget{sourcepointer-2}{}\fi
\noindent\textit{Source reference.} Original remark 0GBG: sources/stacks-chow/tags.txt, source lines 20371--20371. Prerequisite pointer invoked at proof line(s) 4376; it names a standard fact of a sibling chapter that the counted argument uses, and every step of the counted proof built on it is retained verbatim in the delivered body.\par\medskip
\ifdefined\hypertarget\hypertarget{sourcepointer-3}{}\fi
\noindent\textit{Source reference.} Original reference 0GAM: sources/stacks-restricted/pinned-original-source.tex, source lines 1043--1043. Prerequisite pointer invoked at proof line(s) 4390; it names a standard fact of the same chapter that the counted argument uses, and every step of the counted proof built on it is retained verbatim in the delivered body.\par\medskip
\ifdefined\hypertarget\hypertarget{sourcepointer-4}{}\fi
\noindent\textit{Source reference.} Original reference 0GCI: sources/stacks-restricted/pinned-original-source.tex, source lines 4214--4214. Prerequisite pointer invoked at proof line(s) 4380; it names a standard fact of the same chapter that the counted argument uses, and every step of the counted proof built on it is retained verbatim in the delivered body.\par\medskip
\begin{lemma}
\label{lemma-base-change-rig-smooth-continuous}
Consider the properties $P(\varphi)=$``$\varphi$ is rig-smooth''
and $Q(\varphi)$=``$\varphi$ is adic'' on arrows of $\textit{WAdm}^{Noeth}$.
Then $P$ is stable under base change by $Q$ as defined in
Formal Spaces, Remark
\ref{formal-spaces-remark-base-change-variant-variant-Noetherian}.
\end{lemma}

\begin{proof}
The statement makes sense by Lemma \ref{lemma-rig-smooth-continuous}.
To see that it is true assume we have morphisms
$B \to A$ and $B \to C$ in $\textit{WAdm}^{Noeth}$
and that $B \to A$ is rig-smooth and $B \to C$ is adic
(Formal Spaces, Definition
\ref{formal-spaces-definition-adic-homomorphism}).
Then we can choose an ideal of definition $I \subset B$
such that the topology on $A$ and $C$ is the $I$-adic topology.
In this situation it follows immediately that
$A \widehat{\otimes}_B C$ is rig-smooth over $(C, IC)$ by
Lemma \ref{lemma-base-change-rig-smooth-homomorphism}.
\end{proof}

\end{document}
